
\begin{Aside}

% \subsubsection{\texorpdfstring{2) How to derive your working formula
% from \emph{your}
% framework}{2) How to derive your working formula from your framework}}


Recall from \eqref{eq:identify-alpha-trace}:

\begin{equation}
\bm{\alpha}_{\psi}=\bm{b}_{\Omega}
=
\ShiftB\,\bm{\gamma}+\bm{\varsigma}_{\gamma}\,\varkappa
\qquad\text{and}\qquad
\bm{\alpha}_{r}=-\bm{\gamma}.
% \tag{\ref{eq:identify-alpha-trace}}
\end{equation}

The identification \(\bm\alpha_r=-\bm\gamma\) says: \emph{the
reconstruction injects a rigid plane‑section mode whose amplitude is
\(-\bm\gamma\).}

All that remains is to express \(\bm\gamma\) in terms of the Saint‑Venant
parameters \(\IesanVector\).
That comes from the first boxed equation:

\[\bm{b}_{\Section}=\ShiftB\,\bm\gamma+\bm{\varsigma}_{\gamma}\,\varkappa
\quad\Longrightarrow\quad
\boxed{
\bm\gamma=\ShiftB^{-1}\bigl(\bm{b}_{\Section}-\bm{\varsigma}_{\gamma}\,\varkappa\bigr).
}\]

Putting the two boxes together gives the general ``shear‑gauge rigid
rotation'' formula:

\[
\boxed{
q_{S,\theta\perp}
=
\FrameBasis\times\,\ShiftB^{-1}\bigl(\bm{b}_{\Section}-\bm{\varsigma}_{\gamma}\,\varkappa\bigr).
}
\]

% \subsubsection{\texorpdfstring{3) What about torsion/shear coupling?
% (This is where your current code is \emph{not} fully
% general)}{3) What about torsion/shear coupling? (This is where your current code is not fully general)}}

If your trace allows torsion--shear coupling, you can't treat
\(\Twist\) as ``known independently'' unless you also use your other
amplitude identity (same equation block):

\[\alpha_{\varphi}=\tilde\varkappa
=
(\ShiftB^{\top}\bm{\rho}_{\psi}+ \bm{\varsigma}_{\varkappa})\cdot\bm{\gamma}
+(\bm{\rho}_{\psi}\cdot\bm{\varsigma}_{\gamma}+\lambda_{\varkappa})\varkappa.\]

In other words, for the fully coupled case the ``known data from \(\IesanVector\)''
are \((\bm{b}_{\Section},\tilde\varkappa)\), and you must solve the
\textbf{3×3 coupled system} for \((\gamma,\varkappa)\) before forming
\(q_{S,\theta\perp}=\FrameBasis\times\bm\gamma\).

\begin{quote}
In the fully general trace family with couplings, you must first solve
the coupled
\((\bm{b}_{\Section},\tilde\varkappa)\mapsto(\bm\gamma,\varkappa)\) system;
then \(q_{S,\theta\perp}=\FrameBasis\times\,\gamma\).
\end{quote}


\textbf{B. Compute \(\bm\gamma\) (and \(\varkappa\))
from \(p\) using the \textbf{full} trace matrix
blocks}

From your framework \eqref{eq:identify-alpha-trace}, the Iesan amplitudes

\[
\tilde\varkappa:= a_\vartheta+c_\vartheta \in\mathbb R,
\qquad
\]

relate to the beam strains \((\bm\gamma,\varkappa)\) by the
\textbf{trace‑matrix‑dependent} linear system

\[
\boxed{
\begin{aligned}
\bm{b}_{\Section} &= \ShiftB\,\bm\gamma + \bm{\varsigma}_{\gamma}\,\varkappa,\\[2pt]
\tilde\varkappa&= (\ShiftB^{\top}\bm{\rho}_{\psi}+ \bm{\varsigma}_{\varkappa})\cdot\bm\gamma
+(\bm{\rho}_{\psi}\cdot\bm{\varsigma}_{\gamma}+\lambda_{\varkappa})\,\varkappa.
\end{aligned}
}
\tag{1}
\]

All quantities
\(\ShiftB,\bm{\varsigma}_{\gamma},\bm{\varsigma}_{\varkappa},\lambda_{\varkappa},\bm{\rho}_{\psi}\)
are functions of the trace matrix \(\mathbf{T}_0^{-1}\) via the block
formulas \eqref{eq:Tinv-from-T} and the definitions around \eqref{eq:T0-inv}.

Solve (1) as follows.

Define

\[
\bm g := \ShiftB^{\top}\bm{\rho}_{\psi}+ \bm{\varsigma}_{\varkappa}\in\mathbb R^2,
\qquad
\beta := \bm{\rho}_{\psi}\cdot\bm{\varsigma}_{\gamma}+\lambda_{\varkappa}\in\mathbb R.
\]

Then

\[
\bm\gamma = \ShiftB^{-1}(\bm{b}_{\Section}-\bm{\varsigma}_{\gamma}\varkappa),
\]

and substituting into the \(\tilde\varkappa\) equation gives

\[
(\beta-\bm g\cdot\ShiftB^{-1}\bm{\varsigma}_{\gamma})\,\varkappa
=
\tilde\varkappa-\bm g\cdot\ShiftB^{-1}\bm b.
\]

Hence

\[\boxed{
\varkappa
=
\frac{\tilde\varkappa-\bm g\cdot\ShiftB^{-1}\bm b}{\beta-\bm g\cdot\ShiftB^{-1}\bm{\varsigma}_{\gamma}},
\qquad
\bm\gamma
=
\ShiftB^{-1}\bigl(\bm b-\bm{\varsigma}_{\gamma}\,\varkappa\bigr).
}
\tag{2}
\]
\end{Aside}